\documentclass[11pt]{article}
\usepackage[margin=1in]{geometry}
\usepackage[T1]{fontenc}
\usepackage{amsmath}
\usepackage{booktabs}
\usepackage{hyperref}

\title{CADloop: A Research Scaffold for Parametric CAD Agent Harnesses}
\author{Authors to be determined}
\date{Draft -- no experimental claims}

\begin{document}
\maketitle

\begin{abstract}
This document is a writing scaffold for a planned study of agent harnesses
for parametric three-dimensional CAD construction and editing. The research
scope, method, benchmarks, and evaluation protocol remain to be finalized.
No real CAD benchmark results or state-of-the-art claims are reported here.
\end{abstract}

\section{Introduction}
\textbf{TODO:} Motivate the specific modeling or editing problem. Distinguish
executable geometry, satisfaction of design constraints, and editability.
Write contributions only after their supporting experiments are complete.

\section{Related Work}
\textbf{TODO:} Separate CAD generation and representation methods, datasets
and benchmarks, general agent harnesses, and domain-specific CAD agents.
DeepCAD provides a starting point for CAD sequence literature
\cite{wu2021deepcad}; Fusion 360 Gallery provides human design sequences and
a programmatic reconstruction setting \cite{willis2021fusion}.
These sources are candidates for protocol study, rather than evidence of
completed harness reproductions.

\section{Problem Setup}
\textbf{TODO:} Define task inputs, initial CAD state, actions, permitted
observations, artifact format, units, tolerances, budgets, and termination.
Keep evaluator-only ground truth unavailable to the agent.

\section{Method}
\textbf{TODO:} Describe the implemented harness and its state, feedback,
verification, and recovery rules. Account for the cost of each component.

\section{Experimental Setup}
\textbf{TODO:} Freeze source commits, CAD kernel, model versions, tools,
prompts, data manifests, splits, budgets, seeds, and evaluator versions.
Label official reproductions, CAD adaptations, and local protocols clearly.

\section{Results}
\textbf{TODO:} Add audited results from real CAD experiments, including
sample counts, all failure outcomes, uncertainty, and cost curves. Keep
different task protocols in separate comparisons. Mock smoke outputs must
not be inserted as paper results.

\section{Ablations and Failure Analysis}
\textbf{TODO:} Test each actual component with the same model and budget.
Report negative results and representative failures from frozen task sets.

\section{Limitations}
\textbf{TODO:} State supported CAD representations and operations,
dataset and license constraints, potential contamination, numerical
tolerances, provider nondeterminism, and limits on generalization.

\section{Conclusion}
\textbf{TODO:} Summarize only claims supported by completed experiments.

\bibliographystyle{plain}
\bibliography{references}

\appendix
\section{Reproducibility Materials}
\textbf{TODO:} Include prompts, tool schemas, environment details, complete
hyperparameters, failure definitions, and trace-to-table provenance.

\end{document}
